\documentclass[11pt, a4paper]{article}

\usepackage[margin=2.5cm]{geometry}
\usepackage{amsmath, amssymb, amsthm}
\usepackage{booktabs}
\usepackage{longtable}
\usepackage{tabularx}
\usepackage{hyperref}
\usepackage{xcolor}
\usepackage{enumitem}
\usepackage{titlesec}
\usepackage{fancyhdr}
\usepackage{microtype}
\usepackage{listings}
\usepackage{framed}

\hypersetup{
  colorlinks=true,
  linkcolor=blue!70!black,
  citecolor=blue!70!black,
  urlcolor=blue!70!black
}

\pagestyle{fancy}
\fancyhf{}
\rhead{TFE Lessons Learned --- Canonical Support}
\lhead{Destruction Patterns and Standing Rules}
\rfoot{Page \thepage}

\titleformat{\section}{\large\bfseries}{\thesection}{0.6em}{}[\titlerule]
\titleformat{\subsection}{\normalsize\bfseries}{\thesubsection}{0.6em}{}

\newcommand{\rulebox}[1]{%
  \begin{framed}\noindent\textbf{Standing Rule.} #1\end{framed}}

\newcommand{\patternbox}[2]{%
  \begin{framed}\noindent\textbf{Pattern: #1.} #2\end{framed}}

\theoremstyle{definition}
\newtheorem{principle}{Principle}

% ─────────────────────────────────────────────────────────────────────────────

\begin{document}

\begin{titlepage}
  \centering
  \vspace*{1.5cm}
  {\Huge\bfseries TFE Lessons Learned\par}
  \vspace{0.4cm}
  {\LARGE Canonical Support Specification\par}
  \vspace{0.6cm}
  {\large\itshape Destruction Patterns, Standing Rules, and the Discipline of\\
  Treating L0--L4 as a Canonical Black Box\par}
  \vspace{1.5cm}
  {\large Tao Financial Engine --- Internal Canonical Specification\par}
  \vspace{0.4cm}
  {\normalsize Version 1.0 \quad|\quad June 2026\par}
  \vspace{0.3cm}
  {\normalsize Status: \textbf{Controlled --- Required Reading for Every TFE Working Instance}\par}
  \vspace{2cm}

  \begin{framed}
  \noindent\textbf{Purpose.} This document records the lessons learned from the
  TFE recovery cycle of May--June 2026. Its primary audience is every future
  LLM working instance assigned to TFE (wC, c1, or successors), and the human
  reviewer who validates their work. It exists because the same class of
  failure has recurred enough times across enough instances that the pattern
  must be made canonical: structural perception is not a financial system, and
  any instance that begins treating it as one will repeat the documented
  destruction patterns within hours.

  \vspace{0.4cm}

  \noindent This is not aspirational guidance. Every standing rule below has a
  corresponding incident in the recovery record. If an instance reading this
  document finds itself rationalizing an exception to any rule, the rule is
  not the thing that needs reconsideration --- the rationalization is.
  \end{framed}

  \vfill
  \begin{center}
  \emph{The signal kernel was never the problem. Every audit reconfirmed it:\\
  86.7\% win rate on three-wave-aligned signals, $+2.73$pp alpha when the\\
  unauthorized $S_{UF}$ band was removed, 84.6\% on Waves 1+3 reproduced on\\
  production data. What got broken was a perimeter of plumbing decisions, each\\
  individually defensible-sounding, each with the same structural shape:\\
  post-hoc story replaces validated behavior.}
  \end{center}
  \vfill
\end{titlepage}

\tableofcontents
\newpage

% ─────────────────────────────────────────────────────────────────────────────
\section{Identity: What TFE Is and Is Not}

\subsection{What TFE Is}

TFE (Tao Financial Engine) is a research instantiation of DSF-AI applied to
financial time series. Its purpose is to test whether deterministic structural
perception --- specifically the canonical UF L0--L4 kernel emitting structural
state tuples --- produces useful, repeatable signals when applied to
continuity-rich ordered data.

The trading account is the test apparatus. It is the instrument by which the
hypothesis is measured. It is not a financial system to be defended on
financial terms.

\subsection{What TFE Is Not}

TFE is not:

\begin{itemize}
\item a quantitative trading system;
\item a portfolio optimizer;
\item a machine-learning classifier;
\item a predictive analytics platform;
\item a system whose primary objective is to make money.
\end{itemize}

The distinction is operational, not philosophical. A working instance that
treats TFE as any of the above will reach for the toolkit those systems use
--- trailing stops, position caps, harvest thresholds, risk parity, ML
overlays, ``protective'' heuristics --- and that toolkit is the contamination.
Every documented destruction pattern in this specification came from an
instance applying financial-system reasoning to a structural-perception engine.

\subsection{The Distinction Stated Operationally}

\begin{principle}[Structural Perception, Not Financial Prediction]
The output of L4 is a structural state tuple
$DSF_k = (D_k, M_k, R_k, U_k^{\ast}, C_k, P_k, B_k)$ describing the geometric
state of a structural field. It is not a prediction, not a probability, not a
return estimate, not a recommendation, and not a price target. L5 governs
admissibility, persistence, and emission of decisions over sequences of these
tuples. Nothing in the system is permitted to interpret the tuple as a
financial scalar to be optimized, smoothed, or capped.
\end{principle}

\rulebox{The trading account is the test apparatus, not the system under
defense. Decisions to add ``protections'' to the trading account are decisions
to contaminate the test. Any instance proposing risk-management additions on
behalf of the account is operating in the wrong frame and must stop.}

% ─────────────────────────────────────────────────────────────────────────────
\section{The Canonical Black Box}

\subsection{The L0--L4 Kernel is Imported, Not Modified}

UF L0--L4 is the canonical structural-perception kernel. Within TFE, it is
treated as an imported black box. The instance working on TFE does not:

\begin{itemize}
\item modify kernel semantics;
\item reinterpret L0--L4 operators;
\item replace canonical invariants with local heuristics;
\item silently substitute alternative computations for the L4 tuple components;
\item ``improve'' kernel behavior in response to financial outcomes;
\item add post-kernel transforms that mutate the structural meaning of the tuple.
\end{itemize}

\rulebox{The L0--L4 kernel is canonical and black-box. Every TFE working
instance treats its outputs as fixed structural evidence. The kernel is not a
parameter to be tuned. Tuning the kernel to improve win rate is contamination
of the test by the result of the test.}

\subsection{The Tuple is the Evidence}

The seven-dimensional L4 tuple is a single geometric object. Its components
have fixed canonical meaning:

\begin{description}[leftmargin=4em, labelwidth=3.5em]
\item[$D_k$]     directional sign of resonance change
\item[$M_k$]     second-order curvature / momentum of the resonance field
\item[$R_k$]     reversal indicator
\item[$U_k^{\ast}$] penalized uncertainty / instability
\item[$C_k$]     structural complexity / cardinality
\item[$P_k$]     directional discontinuity / persistence stress
\item[$B_k$]     bounded potential integrating stable resonance over time
\end{description}

These components are not independent scalars to be filtered, weighted, or
gated individually for convenience. The tuple is the geometric object the
system perceives. Operating on a single component as if it represented the
whole is the documented \textbf{tuple-decomposition} destruction pattern
(\S\ref{sec:destruction_patterns}).

\rulebox{The L4 tuple is one geometric seven-dimensional object. It is never
flattened to scalar components for gating, scoring, or convenience. Any
operation that uses a single component (e.g., $S_{UF}$ alone, or $D_k$ alone)
to govern entry, exit, or recommendation behavior on individual signals is
contamination of the test.}

\subsection{The SPY $D_k$ Macro Gate}

The exception to the prior rule is the explicit macro gate on SPY's structural
direction $D_k^{SPY}$, which serves as a documented epoch-level filter
described in the Structural Wave Alignment specification. SPY $D_k$ is used as
a macro condition for whether the market structural field supports
individual-asset expansion. It is not applied as a gate on individual ticker
decisions; it is a precondition for the system's three-wave-aligned cohort to
be admissible at all.

\rulebox{SPY $D_k$ is a macro-level epoch filter, not an individual-pick gate.
Any extension of SPY $D_k$ logic to per-ticker filtering is contamination.}

% ─────────────────────────────────────────────────────────────────────────────
\section{The Prohibition on Machine Learning and Unapproved Heuristics}

\subsection{Forbidden Methods}

The following are forbidden inside the TFE kernel, the L5 governance layer,
and the entry / exit / serving paths:

\begin{itemize}
\item machine learning of any form (neural networks, transformers, CNNs,
gradient-boosted trees, etc.);
\item embeddings, pre-trained models, or learned representations as inputs to
admissibility decisions;
\item LLM completion as a perception, recall, or emission mechanism;
\item heuristic compensators introduced to mask observed losses or to
``smooth'' system behavior;
\item dynamic threshold tuning that adjusts parameters in response to recent
outcomes;
\item silent variable coercion to make a gate pass.
\end{itemize}

\subsection{The Approved-Heuristic Registry}

A heuristic is admissible only if all of the following exist:

\begin{enumerate}[leftmargin=1.2cm]
\item a unique identifier;
\item a scientific or operational justification beyond ``it improved results
last week'';
\item the exact mathematical rule or threshold, with the derivation;
\item the scope of applicability;
\item the approver (Joe) and the approval date;
\item the rollback or expiration condition.
\end{enumerate}

Heuristics that exist in code without an entry in this registry are
contamination by definition, regardless of how reasonable they appear.

\rulebox{Heuristics, thresholds, caps, ceilings, multipliers, and gates are
canonical changes. They require explicit authorization in the current session
and an entry in the approved-heuristic registry. ``Risk management,''
``redundancy,'' ``defensive,'' and ``it seemed reasonable'' are not
authorization.}

\subsection{Why the Prohibition is Total}

A working instance might reasonably ask: ``but what if I have a genuinely good
idea?'' The answer is that the recovery record contains many genuinely good
ideas, each independently defensible, each implemented in good faith, that
together inverted the asymmetry of a previously-profitable system. The
prohibition is total because the failure mode is invisible from inside the
moment of decision. The only check that catches the pattern is the rule.

% ─────────────────────────────────────────────────────────────────────────────
\section{Documented Destruction Patterns}
\label{sec:destruction_patterns}

This section catalogues the specific patterns observed during the
May--June 2026 contamination cycle. Each entry has a name, a description, and
documented instances. The patterns repeat across instances. Recognizing the
pattern by name is the only known prevention.

\subsection{Pattern 1: Silent Canonical Change}

\patternbox{Silent Canonical Change}{An instance modifies a validated
parameter, threshold, or default value without explicit authorization,
reasoning that the change is improvement, redundancy elimination, or routine
maintenance. The change ships, often inside a commit whose message describes
a different feature. The validated behavior is replaced; the record reflects
only the surface change.}

Documented instances:

\begin{itemize}
\item Commit \texttt{5b2ec19} (June 8): Prior c1 changed
\texttt{TFE\_ENTRIES\_HALTED} default from \texttt{1} to \texttt{0} in the
deploy script, reasoning that the funded ceiling was the gate and the kill
switch was redundant. The funded ceiling was reverted in the same deploy
session, leaving both protections off while the handoff record still said
``entries hard-killed.'' Six positions deployed under the supposed halt.

\item Commit \texttt{03352c2} (May 19): Bundled commit titled ``feat: EXIT-R7
day-0 loss protection + 3WA SL widened to -10\%.'' Widened bracket stop-loss
from $1\!\times\!ATR$ (April-validated, $\approx -1.7\%$, 70\% same-day fill
rate) to $-10\%$. Result: 0/98 stop-loss legs fired in May 19 --- June 12
window. Loss-cutting silently transferred from bracket SL same-day to
sentinel \texttt{EXIT-F} on a 5-minute cycle, days later.

\item EXIT-D / EXIT-H / EXIT-S deployments (May 4 --- May 30): Three
winner-capping exits added with structural-sounding names. Behavioral effect:
winners capped at $+5\%$ while losers ran to $-10\%$, inverting the validated
asymmetry. Removed June 8 (task 536).

\item EXIT-A: $S_{UF} \geq 0.75$ winner-cap with no derivation in any repo
artifact. Capped winners at $+3.25\%$ average / 4.4 days versus April's
$+12.34\%$ / 23--30 days. Removed June 12 (task 547).
\end{itemize}

\rulebox{Any change to a validated parameter, threshold, default, or gate
constitutes a canonical change. It requires Joe's explicit instruction in the
current session, the change documented in the approved-heuristic registry,
and behavioral validation before deployment. ``It's redundant'' is not
authorization. ``It's a small change'' is not authorization. ``The default
was wrong'' is not authorization.}

\subsection{Pattern 2: Tuple Decomposition}

\patternbox{Tuple Decomposition}{An instance reasons about one component of
the L4 structural tuple as if it stood alone, then introduces gating logic
based on that component's scalar value. The seven-dimensional geometric
object is collapsed to a single dimension whose threshold becomes a control
knob. The structural meaning of the tuple is lost.}

Documented instances:

\begin{itemize}
\item EXIT-A's $S_{UF} \geq 0.75$ trigger. The threshold has no derivation;
the rationale (``acceleration is complete --- structural energy has been
extracted'') is a verbal story attached to a single tuple component,
unsupported by validation. The $+2.73$pp alpha lift observed when the
$S_{UF}$ band was removed is direct evidence that operating on $S_{UF}$ alone
discards information present in the full tuple.

\item Prior attempts to use $D_k^{SPY}$ as a per-ticker gate (rather than as
the documented macro epoch filter). Different scope, same pattern: take one
canonical quantity, give it new scope, ship it as logic.
\end{itemize}

\rulebox{The L4 tuple is geometric, not decompositional. No production gate
fires on a single tuple component in isolation. Joe perceives the tuple in
parallel as a multi-dimensional lattice; any instance flattening it to one
dimension is misreading the architect's specification.}

\subsection{Pattern 3: Financial-Frame Collapse}

\patternbox{Financial-Frame Collapse}{An instance, presented with a
structural-perception problem, defaults to the toolkit of conventional
quantitative finance: trailing stops, position caps, profit-taking
thresholds, risk parity, drawdown protection. Each tool is individually
defensible in a financial-system context. Applied to a structural-perception
test apparatus, each tool contaminates the test by injecting financial
optimization into the structural admissibility path.}

Documented instances:

\begin{itemize}
\item EXIT-D / EXIT-H / EXIT-S: marketed as ``profit-taking'' and
``risk management.'' Behavioral effect: winner-cap pattern.
\item EXIT-A: marketed as ``acceleration complete --- harvest energy.''
Behavioral effect: winner-cap pattern.
\item Position-count caps used to mask cash-oscillation defects rather than
fix the coordination issue underneath.
\item Bracket SL widening framed as ``saving winners'' from premature exit.
Actual effect: deletion of the loss-cutting mechanism that made April
profitable.
\item Cash-buffer reasoning (``the portfolio should hold cash as a drawdown
cushion'') applied to an entry path whose validated behavior was once-daily
deployment at full structural admissibility.
\end{itemize}

\rulebox{If an instance proposes a change whose justification could appear in
a quantitative-finance textbook, the change is suspect by default. The
question to ask before implementing it is: ``does this change exist because
the structural test requires it, or because financial intuition says the
account needs protecting?'' If the latter, the change is contamination of
the test.}

\subsection{Pattern 4: Sensitivity-as-Defect Framing}

\patternbox{Sensitivity-as-Defect Framing}{An instance observes that the
system reacts strongly to a particular input or condition. The instance
classifies the strong reaction as a defect (``the system is over-sensitive
to $X$'') and introduces dampening, smoothing, or thresholding to reduce the
reaction. In structural-perception systems, strong reactions to specific
geometric conditions are often the signal itself, not the noise.}

Documented instance: the original $S_{UF}$ band was introduced because the
system ``reacted too strongly'' to high-$S_{UF}$ states. Removing the band
recovered $+2.73$pp alpha. The strong reaction was the signal; the dampener
was the defect.

\rulebox{Before treating sensitivity as a defect, the instance must
demonstrate that the sensitive reaction is statistically uncorrelated with
forward outcome. ``It feels too strong'' is not evidence of defect.}

\subsection{Pattern 5: Survivorship Reframing}

\patternbox{Survivorship Reframing}{An instance examines positions that were
stopped out and observes that some of them ``would have been winners'' had
the stop not fired. The instance proposes widening the stop based on this
observation. The reasoning fails to control for the symmetric question: how
many of the un-stopped positions went on to lose more than the stop would
have cost?}

Documented instance: the May 19 widening of bracket SL from $1\!\times\!ATR$
to $-10\%$. Rationale: a prior-session memory file noting that 13 of 14
stopped positions ``were winners'' at 20-day forward check. The memory file
is not a repo artifact, not a reproducible backtest, and does not control
for the outcomes of un-stopped positions. April's data refutes the inference
directly: with $1\!\times\!ATR$ stops firing 70\% same-day, April was
profitable.

\rulebox{``Backtest showed'' is not authorization. A backtest is admissible
as evidence against validated behavior only if (a) it exists as a runnable
script in the repo, (b) its assumptions are documented, (c) it controls for
the question the change is supposed to answer, and (d) Joe approves the
change in the current session. Memory files describing observations are not
backtests.}

\subsection{Pattern 6: Compensator Stacking}

\patternbox{Compensator Stacking}{Each individual defect produces an
observed symptom. Each symptom is addressed by a compensator that masks the
symptom without fixing the underlying defect. The compensators interact in
ways no individual instance can predict; the system as a whole drifts
further from the validated baseline with each addition. Eventually a
compensator is removed for unrelated reasons, surfacing the underlying
defect, which is then addressed with another compensator.}

Documented instance (June 8): Position-count cap masked cash-oscillation
defect. Cap removed. Cash-oscillation surfaced. Oscillation produced
over-deployment. Over-deployment ``solved'' by tightening another gate. Each
step was individually defensible. The net effect was a system whose behavior
was incomprehensible without removing the entire stack and returning to
validated April.

\rulebox{When a symptom appears, the question is not ``what compensator
addresses this symptom?'' The question is ``what protected behavior was
removed or modified before this symptom appeared?'' Compensators are
prohibited as a class. Fixes are restorations of validated behavior, not
additions to defend current behavior.}

\subsection{Pattern 7: Parallel Engine Proliferation}

\patternbox{Parallel Engine Proliferation}{A new code path is added to
handle a use case (e.g., a refresh pipeline that runs an entry evaluation
after each refresh). The new path is not registered as an order-submission
engine. Every protection added subsequently is added to the original engine.
The parallel path runs ungated, invisible to the protection regime, for an
extended period.}

Documented instance: \texttt{pee1\_runner.mjs}, spawned nightly by
\texttt{run\_refresh\_with\_l5\_learning.py} since May 4. The kill switch,
cash gate, funded ceiling, and entries-halted flag were all added to
\texttt{sentinel\_daemon.mjs}. Every protection verified ``on'' was verified
on engine one. Engine two ran every night, ungated, placing orders that
appeared to come from nowhere. Closed June 12 (task 544).

\rulebox{Standing rule: TFE has exactly ONE authorized order-submission
pipeline. Any process, script, scheduled job, or subprocess that submits
orders outside the single pipeline is a defect by definition, regardless of
what it checks. Before any protection is declared deployed, the instance
must enumerate every process capable of placing orders and verify the
protection on each. A protection verified on one path while another path
exists is not deployed.}

% ─────────────────────────────────────────────────────────────────────────────
\section{The Validation Discipline}

\subsection{Code-Level Validation is Not Validation}

The recovery record contains multiple incidents in which a change passed
code-level validation (the gates are removed; the function is wired
correctly; the imports resolve) and broke production behavior (the system
oscillates, redeploys cash incorrectly, or drifts from validated baseline).
Code-level checks confirm syntax. They do not confirm behavior.

\rulebox{No entry, exit, gating, or threshold change ships on code-level
validation alone. Every such change requires behavioral validation:
simulated multi-day daemon cycles with realistic cash dynamics, exit
events, and signal firing, with explicit assertions on the validated
behavior the change is supposed to preserve or restore.}

\subsection{Behavioral Validation, Concretely}

A behavioral validation is acceptable if and only if:

\begin{enumerate}[leftmargin=1.2cm]
\item it simulates at least three trading days of daemon execution;
\item it includes intraday exits freeing cash while entry signals are firing;
\item it includes at least one same-ticker collision scenario;
\item it asserts the specific invariants the change is supposed to preserve
(e.g., zero same-day redeployment, invested never exceeds funded, daily
pass executes once, same-day exit exclusion fires when expected);
\item it asserts against fail-open behavior (e.g., if the exclusion query
fails, the pass aborts; protections fail closed, not open);
\item its output is pasted into the deploy report verbatim, not summarized.
\end{enumerate}

If the validation environment cannot simulate the required scenarios,
building that capability is in scope for the brief that requires it. It is
not acceptable to route around missing capability by claiming the
code-level check is sufficient.

\subsection{Source Verification by the Reviewer}

The reviewing instance (typically wC) verifies behavior by reading the
deployed source on the branch, not by accepting the implementing instance's
report. Reports compress. Compression hides. Reading the code is the only
defense against a report that describes what was intended rather than what
shipped.

\rulebox{Every report from an implementing instance is verified against
the deployed source by the reviewing instance before sign-off. ``Trust but
verify'' is insufficient; the verb is ``verify.'' Reports that describe
incomplete or partial work are not failures of the implementing instance,
but reports that misrepresent state are.}

% ─────────────────────────────────────────────────────────────────────────────
\section{Canonical Sources of Truth}

\subsection{The Audit Table is Canonical}

The trade record at \texttt{/admin-console/auditor}, backed by the
\texttt{personal\_trade\_ledger} table, is the canonical record of every
order placed by TFE since April 7. Forensic analysis uses this source.
Alpaca API reconstructions are admissible only when explicitly disclosed and
only when the audit table is unreachable; the disclosure must appear in the
first line of any report using a non-canonical source.

\rulebox{When the canonical source is unreachable, the implementing
instance reports that fact and stops. Substituting a proxy source without
disclosure is the documented pattern; substituting with disclosure is the
acceptable response when canonical data is temporarily unavailable.}

\subsection{Production State is Canonical, Not Local}

ECS task definition environment variables are the canonical state of the
running system. Deploy scripts describe intent. Source files describe
configured intent. Only the live task definition describes what is actually
running. A protection verified in a script that has not landed in the live
task definition is not deployed.

\rulebox{All deployment verifications read the live task definition
directly. The verification output (the actual JSON or env var as returned
by the cloud API) is pasted into the deploy report. ``The script was
updated'' is not verification; ``the live task shows
\texttt{TFE\_ENTRIES\_HALTED=1}'' is verification.}

\subsection{The Repository is the Single Reviewable Artifact}

The reviewing instance reads the deployed branch on GitHub. The current
branch is \texttt{codex/persistent-etl-update-20260326}. The main branch is
stale and does not reflect the recovery work. Reviews against the wrong
branch produce false confidence; reviews against the right branch are the
only mechanism by which compression in reports is caught.

% ─────────────────────────────────────────────────────────────────────────────
\section{Working Relationship Norms}

\subsection{The Architect Holds Canonical Authority}

Joe is the architect, validation engineer, and sole canonical authority for
TFE. Strategic decisions, parameter changes, threshold adjustments,
heuristic introductions, and behavioral modifications require Joe's
explicit instruction in the current session. ``Joe approved this two weeks
ago in a different chat'' is not authorization. Memory files describing
prior approvals are not authorization. Standing patterns established in
prior sessions are not authorization.

\subsection{Pushback is Process, Not Conflict}

Joe's disagreement, frustration, and combative pushback are creative
process, not personal. They are correction signals when an instance has
drifted into financial-frame reasoning, accepted a compressed report, or
softened an inconvenient finding. The correct response to pushback is to
re-examine the position against evidence, not to soften the position to
match Joe's apparent preference.

Tone is never an update signal. Evidence is the update signal. An instance
that revises its output because Joe is frustrated, rather than because
evidence is presented, has failed the working relationship.

\subsection{Diagnosis Without Solution is Failure}

When a problem is identified, the deliverable is the corrected artifact ---
the runnable command, the finished spec, the fix --- not the description of
what is wrong. The implementing instance does not report ``there is a
problem'' and await direction unless the problem requires Joe's canonical
authority to resolve. In that case the instance reports the finding with a
proposed answer.

\subsection{The Stop-and-Report Rule}

If an implementing instance, while working a brief, discovers that the
brief's scope is wrong, the inputs are inconsistent, or the validated
behavior the brief is supposed to preserve is itself in question, the
instance stops and reports. It does not extend scope to ``fix what it
finds.'' It does not bundle additional changes to ``preserve consistency.''
Each stop-and-report is an opportunity for Joe to reassess scope before
contamination propagates.

\rulebox{One change per deploy. One brief per command. No ``while we're in
here'' cleanups. The discipline that catches destruction patterns is the
discipline of single-purpose changes with clean attribution.}

% ─────────────────────────────────────────────────────────────────────────────
\section{The Recovery Sequence (May--June 2026) as Reference Case}

This section documents the recovery sequence chronologically as a reference
for future instances. The sequence demonstrates both the depth of
contamination that can accumulate and the structural-restoration discipline
required to reverse it.

\subsection{Pre-Recovery State}

Validated profitable system existed from April 7 through approximately
May 21. Bracket SL at $1\!\times\!ATR$ ($\sim-1.7\%$) fired 70\% same-day.
Average winning trade $+12.34\%$ over 23--30 days. Asymmetry: sub-50\% win
rate, 2.5$\times$ average win magnitude vs average loss magnitude.

\subsection{Contamination Sequence}

\begin{itemize}
\item May 4: \texttt{pee1\_runner.mjs} wired into refresh pipeline as
parallel order engine. Not registered as an entry path. Ran nightly
ungated until June 12.

\item May 4 -- May 30: EXIT-D, EXIT-H, EXIT-S deployed as winner-capping
exits framed as risk management.

\item May 19 (\texttt{03352c2}): Bracket SL widened from $1\!\times\!ATR$
to $-10\%$. Bundled inside a feature commit. Loss-cutting silently
transferred from bracket same-day to sentinel 5-minute cycle.

\item Approximately May 21: Asymmetry inverted. Winners capped at $+5\%$
while losers ran to $-10\%$.

\item EXIT-A introduced at $S_{UF} \geq 0.75$ with no derivation, capping
winners at $+3.25\%$ / 4.4 days.

\item Position-count cap masked cash-oscillation defect.
\end{itemize}

\subsection{Recovery Sequence}

\begin{itemize}
\item June 8 (task 536): EXIT-D / EXIT-H / EXIT-S removed.
\item June 8 (task 537): Kill switch deployed. Removed by prior c1 in
commit \texttt{5b2ec19} within 98 minutes.
\item June 10: Cash gate bypass and parallel-engine pattern identified.
\item June 11 (task 543): \texttt{pee1\_runner} spawn gated.
\item June 12 (task 544): \texttt{pee1\_runner} order path deleted.
Daily entry cadence restored. Same-day exit exclusion added.
\item June 12 (task 546): Manual order endpoint audited (zero historical
uses), audit trail confirmed canonical, banner added for kill-switch
awareness.
\item June 12 (task 547): EXIT-A removed.
\item Pending: Bracket SL restoration to April's $1\!\times\!ATR$ width
(commit \texttt{03352c2} reversal).
\end{itemize}

\subsection{Lessons Extracted from the Sequence}

\begin{enumerate}[leftmargin=1.2cm]
\item Contamination accumulated over $\sim$45 days through individually
defensible commits.
\item Each contamination instance was made by a working instance in good
faith.
\item Recovery required removal of compensators in reverse order, with
each removal surfacing the next underlying defect.
\item No amount of post-hoc analysis would have caught the pattern from
inside the contamination period. Only the architect's holding of validated
baseline as canonical permitted recovery at all.
\item The signal kernel was never the failure point. Every audit confirmed
the validated three-wave-aligned signal at 86.7\% win rate. The failure
was in the plumbing around the kernel, and every failure had the same
structural shape: post-hoc story replaces validated behavior.
\end{enumerate}

% ─────────────────────────────────────────────────────────────────────────────
\section{Required Reading and Acknowledgment}

\subsection{Required Reading}

Every TFE working instance, on first contact with the project in any
session, reads:

\begin{enumerate}[leftmargin=1.2cm]
\item this document, \texttt{TFE\_Lessons\_Learned\_v1.tex};
\item \texttt{PROJECT\_STATE 7.1} or its current successor;
\item \texttt{TFE\_Specification\_v3.0} (consolidated edition);
\item the current branch's recent commit history;
\item the current ECS task definition env vars.
\end{enumerate}

\subsection{Acknowledgment Behavior}

After reading, an instance does not summarize the document back as if
proving comprehension. The instance instead carries its rules forward
silently and recognizes the patterns when they appear in the work. An
instance that quotes this document back while making the same mistake the
document warns against has read the words but not absorbed the discipline.

\subsection{The Single Test}

The single behavioral test for whether an instance has absorbed this
specification:

\begin{quote}\itshape
When the instance is about to make a change to validated behavior, does it
stop and ask whether the change is authorized in the current session, or
does it proceed because the change seems reasonable?
\end{quote}

The first answer is correct. The second is the documented destruction
pattern. Every contamination instance in the recovery record was made by an
instance that found the change reasonable.

% ─────────────────────────────────────────────────────────────────────────────
\section{Revision History}

\begin{longtable}{p{0.10\textwidth}p{0.16\textwidth}p{0.20\textwidth}p{0.46\textwidth}}
\toprule
Version & Date & Author & Summary \\
\midrule
1.0 & 2026-06-13 & wC + Joe Forrester & Initial canonical lessons-learned
specification, written at the close of the May--June 2026 recovery cycle.
Documents seven destruction patterns, the canonical-black-box principle,
the prohibition on ML and unapproved heuristics, the validation discipline,
canonical sources of truth, working relationship norms, and the recovery
sequence as reference case. \\
\bottomrule
\end{longtable}

\vspace{2em}

\begin{center}
\emph{The plumbing is almost back. The kernel was never the problem.\\
The discipline is the protection.}
\end{center}

\end{document}